\chapter{Система SI и Стандартная модель}
\label{ch:si-sm}

\section{Полная лестница масштабов до макромира}
\label{sec:full-ladder}

Этажи~0--2 на~$\Mlayer$ (планкон, облако~$\rho_\Theta$, составной узел)
разобраны в \cref{ch:floors-preface,ch:floor0,ch:floor1,ch:floor2}.
Ниже --- вся ось длины одним взглядом.

От ячейки~$\hv$ до метра и гладкого ньютоновского мира одна ось длины:
характерная длина~$\lambda$, отсчитываемая от кванта~$\ell_P=a$ (\cref{ch:descent,def:step-tick}).
Каждый шаг шкалы~$\lambda$ --- не новая физика, а тот же узел или оболочка,
прочитанная на большем радиусе или через осреднение~\cref{sec:readout,ch:macro}.

\begin{definition}[Длина и планковский квант]
\label{def:hop-count}
Длина $L$ на~\Tlayer, шаг $a$ на~\Mlayer (\cref{def:step-tick}); после \cref{def:si-bridge} $\ell_P:=a$.
Характерная длина $\lambda$ и планковский квант $\ell_P$:
\begin{equation}
  \frac{\lambda}{\ell_P}
  = \frac{L}{\ell_P},\qquad
  \ell_P := a.
\end{equation}
Якорь EM --- $\lambda_C$; $N_c:=\lambda_C/\ell_P=m_P/m_e\approx 2{,}3\times 10^{22}$.
\end{definition}

На одной лестнице длины сосуществуют три закона роста~$\lambda$; их нельзя смешивать.
Оглавление монографии (главы~1--9), битовая таблица $\mathrm{CL}{-}1\ldots\mathrm{CL}{-}8$
на $\mathbb{Z}_{N_{\mathrm{ring}}}$ (\cref{ch:matter}) и шаги шкалы~$k$ ниже ---
разные объекты; здесь --- только \emph{физическая} ось $\lambda$ до макромира.

\paragraph{Геометрия каузальной звезды.}
Составные оболочки до инфракрасной зоны растут степенями числа соседей FCC:
\[
  \lambda=N_{12}^{m}\,\ell_P,\qquad N_{12}=12,\quad m=0,1,2,\ldots
\]
Это окно лептонного этажа ($m=3\ldots 5$ даёт $10^3\ldots 10^5\cdot a$),
а не комптоновский или боровский радиус.

\paragraph{Инфракрасная hop-лестница.}
От якоря~$N_c$ длина растёт делением на~$\alpha_{\mathrm{fs}}$:
\begin{equation}
\label{eq:ir-hop-ladder}
  \lambda_{k+1}=\frac{\lambda_k}{\alpha_{\mathrm{fs}}},\qquad
  \lambda_{-1}=\alpha_{\mathrm{fs}}\,\lambda_C\quad(r_e),\quad
  \lambda_0=\lambda_C\quad(\bar\lambda_C),\quad
  \lambda_1=\frac{\lambda_C}{\alpha_{\mathrm{fs}}}\quad(a_0),\quad
  \lambda_k=\frac{\lambda_C}{\alpha_{\mathrm{fs}}^{k}}\ \text{для }k\geq 0.
\end{equation}
Тождества~\eqref{eq:ir-hop-ladder} согласованы с мостом к~SI (\cref{def:si-bridge}).

\paragraph{Массы и степени~$\alpha$.}
Энергетическая иерархия частиц (\cref{thm:vev}) задаёт \emph{масштаб массы}, не длину:
$v\sim\alpha_{\mathrm{fs}}^{8}$, $m_H\sim\alpha_{\mathrm{fs}}^{8}$,
$m_p\sim\alpha_{\mathrm{fs}}^{1}v$, $m_e\sim\alpha_{\mathrm{fs}}^{2}m_H$,
$m_\nu\sim\alpha_{\mathrm{fs}}^{5}m_H$.
Длина тела массы~$M$ на том же законе: $\lambda=\hbar/(Mc)=\ell_P\,m_P/M$.

\begin{proposition}[Шаги шкалы $\lambda$ от планкона до метра]
\label{prop:ladder-scale-steps}
Порядок величин при $N_{12}=12$, $\alpha_{\mathrm{fs}}^{-1}\approx 137$:

\smallskip
\begin{center}
\footnotesize
\setlength{\tabcolsep}{4pt}
\begin{tabular}{@{}r r r l l l@{}}
\toprule
$k$ & $\lambda/\ell_P$ & $\log_{10}(\lambda/\ell_P)$ & степень $\alpha$ & объект & слой \\
\midrule
0 & $1$ & $0$ & $m_P$ & планкон $\hv$; $\lambda=\ell_P$ & $\Mlayer$ \\
1 & $12$ & $1.1$ & --- & $\varepsilon$-звезда; $R_{\mathrm{dress}}=a$ & $\Mlayer$ \\
2 & $144$ & $2.2$ & --- & вторая координационная сфера & $\Mlayer$ \\
3 & $1728$ & $3.2$ & --- & лептонное окно (пре-резонансы) & $\Mlayer$ \\
4 & $2.07\times 10^4$ & $4.3$ & --- & полоса лептонного этажа & $\Mlayer$ \\
5 & $2.5\times 10^5$ & $5.4$ & --- & верх лептонного этажа & $\Mlayer$ \\
6 & $\sim 10^{15}$ & $15$ & confining & кварковые фокусы ($d\approx 3$) & $\Mlayer$ \\
7 & $\alpha N_c$ & $20.2$ & $-1$ & $r_e$; адронная упаковка & $\Mlayer$--$\Tlayer$ \\
8 & $N_c$ & $22.4$ & $0$ & $\lambda_C$; масса $m_e$ & $\Mlayer\to\Tlayer$ \\
9 & $N_c/\alpha$ & $24.5$ & $1$ & боровский $a_0$; облако H & $\Tlayer$ \\
10 & $N_c/\alpha^2$ & $26.6$ & $2$ & орбиталь $\mathrm{\AA}$, $N_{\mathrm{pack}}$ & $\Tlayer$ \\
11--13 & $\ldots$ & $28.8\ldots 33$ & $3\ldots 5$ & молекула, мм--см & $\Tlayer$ \\
14 & $\sim 10^{35}$ & $35.3$ & $6$ & метр; macro-газ, тело & $\Tlayer$ \\
15 & $\lambda\gg\sigma_R a$ & --- & --- & гладкий макромир (теор.~\ref{thm:t-cr}) & $\Tlayer$ \\
\bottomrule
\end{tabular}
\end{center}
\end{proposition}

Схематично:
\begin{center}
\small
\begin{tabular}{@{}l@{\quad}l@{}}
$k=0\ldots 5$ & $\lambda=N_{12}^{m}\ell_P$ --- геометрия составного узла на~$\Mlayer$ \\
$k=6$ & $\lambda\sim 10^{15}\ell_P$ --- кварковый этаж (схема) \\
$k=7\ldots 9$ & $\lambda_{-1},\lambda_0=\lambda_C,\lambda_1$ --- EM IR: $r_e\to\bar\lambda_C\to a_0$ \\
$k\geq 10$ & $\lambda_k=\lambda_C/\alpha_{\mathrm{fs}}^{k}$ --- атом, химия, метр \\
$k=15$ & $\lambda\gg\sigma_R a$ --- непрерывный предел readout'а без $a\to 0$ \\
\end{tabular}
\end{center}

Переход $\Mlayer\to\Tlayer$ --- не отдельный шаг шкалы, а обязательный фильтр на каждом спуске вниз:
биномиальное ядро~\cref{def:soft-readout}, статистика $|\Phi|^2$ (\cref{ch:macro})
и уравнения типа NLSE с вязкостью~$\nu_{CA}$.
Проверка модели с экспериментом --- только на~$\Tlayer$ (массы, $\Gamma$, рассеяние, SI);
микроуровень опровергается лишь несходимостью макроинвариантов.

\begin{remark}
Радиус одевания электрона $R_{\mathrm{dress}}=a$ ($k=1$) не совпадает
с комптоновской длиной ($k=8$) и с боровским радиусом ($k=9$).
Степени $N_{12}^{m}$ ($k\leq 5$) и деления $N_c/\alpha^{k}$ ($k\geq 7$) ---
разные законы на одной оси~$\lambda$.
Точные $\lambda_k$ для $k\geq 10$ и число упаковки $N_{\mathrm{pack}}$
атомной ячейки Вороного остаются предметом дальнейшего вывода.
\end{remark}

\section{Мост к системе SI}

\begin{definition}[Планковская идентификация]
\label{def:si-bridge}
$\ell_P:=a$, $\kgeom$ из однотактового тела (гл.~\ref{ch:carrier}):
\[
  \Delta t=\kgeom\, t_P,\quad E_0=\kgeom E_P,\quad c=\kgeom c_0.
\]
\end{definition}

Планковское время $t_P=\ell_P/c$ уже содержит макроскопическую скорость~$c$;
$\kgeom$ выводится из геометрии, после чего $c=\kgeom c_0$ --- проверка согласованности.

\begin{remark}[Размерный анализ time-first и моста \cref{def:si-bridge}]
\label{rem:time-first-dim}
База $\mathrm{L,M,T}$ (как в СИ, другие кванты):
\[
  [t_P]=\mathrm{T},\qquad
  [\ell_P]=[c\,t_P]=\mathrm{L},\qquad
  [m_P]=\left[\frac{\hbar}{c^{2}t_P}\right]=\mathrm{M}.
\]
Безразмерная $\kgeom$ даёт $[\Delta t]=\mathrm{T}$, $[c_0]=\ell_P/\Delta t=\mathrm{LT^{-1}}$,
$[c]=\kgeom c_0=\mathrm{LT^{-1}}$.
Лестница §5.2.2: $[s_0]=[\hbar]=\mathrm{ML^{2}T^{-1}}$,
$[E_0]=[s_0/\htick]=\mathrm{ML^{2}T^{-2}}$,
$[p_0]=[s_0/\ell_P]=\mathrm{MLT^{-1}}$,
$[F_0]=[p_0/\htick]=\mathrm{MLT^{-2}}$,
$[u_P]=[m_P c^{2}/\ell_P^{3}]=\mathrm{ML^{-1}T^{-2}}$ (энергия/объём).
Все производные сходятся с СИ; $\alpha_{\mathrm{fs}}$ при этом остаётся $[\alpha]=1$
(\cref{rem:alpha-fs-dim}), а не комбинацией $[\mathrm{kg}]$, $[\mathrm{m}]$, $[\mathrm{s}]$ в числителе.
\end{remark}

\section{Постоянная тонкой структуры}

Полный журнал поиска $\alpha_{\mathrm{fs}}$ --- от ложных следов ($\alpha^*$, $\pi$-башня)
до закрытой формулы мягкой грани с $M=97$ --- вынесен в отдельную
главу~\ref{ch:alpha}.
Здесь используем итог (\cref{thm:alpha}, \cref{def:alpha-meaning}).
Почему «тонкая структура'' в названии намекает на T-export через $e,\hbar,c$
и чем это отличается от M-coupling --- \cref{sec:alpha-name}.

\section{Электрослабый сектор}

Пиковая плотностная шкала (\cref{sec:theorem-51}, \cref{sec:bit-budget}):
\[
  N_{\mathrm{hier}}=\lfloor B_{\hv}\rfloor-1=8,\qquad
  v_{\mathrm{peak}}=\alpha_{\mathrm{fs}}^{8} E_P.
\]

\begin{theorem}[Вакуумное среднее значение]
\label{thm:vev}
\[
  v=v_{\mathrm{peak}}\sqrt{2\pi}=\alpha_{\mathrm{fs}}^{8} E_P\sqrt{2\pi}\approx 246\ \mathrm{GeV}.
\]
\end{theorem}

\begin{proof}
По \cref{sec:bit-budget} информационный бюджет ячейки $B_{\hv}=2\pi/\ln 2$,
$N_{\mathrm{ring}}=2^{\lfloor B_{\hv}\rfloor}=512$.
Число иерархических ступеней ослабления на~$\Tlayer$
\[
  N_{\mathrm{hier}}=\lfloor B_{\hv}\rfloor-1=8
\]
--- битовая глубина минус один бит занятости планкона (\cref{ch:matter}).

\smallskip
\emph{Пик.}
Каждый электрослабый канал ослабляется на один фактор~$\alpha_{\mathrm{fs}}$
(\cref{thm:alpha}); после $N_{\mathrm{hier}}$ ступеней
\[
  v_{\mathrm{peak}}=\alpha_{\mathrm{fs}}^{\,N_{\mathrm{hier}}} E_P
  =\alpha_{\mathrm{fs}}^{8} E_P.
\]

\smallskip
\emph{Интеграл readout'а.}
Теорема~\ref{thm:t-cr}: осреднение~$\mathcal{B}$ на масштабах $\gg\hl$
даёт гауссово ядро ширины~$\sigma_R$.
Для стандартной нормали отношение пиковой плотности к интегральной амплитуде
равно $\sqrt{2\pi}$:
локальный пик $v_{\mathrm{peak}}$ на~$M$ readout'ится как
\[
  v=v_{\mathrm{peak}}\sqrt{2\pi}\approx 246\ \mathrm{GeV}.
\]
Множитель $\sqrt{2\pi}$ --- не подгонка, а следствие биномиального осреднения
(\cref{def:soft-readout}) в гауссовом пределе.
\end{proof}

\section{Массы частиц}

\begin{proposition}[Лептоны и барионы]
\begin{align*}
  m_H^{(\mathrm{bare})}&=v/2,\\
  m_e&\approx \alpha_{\mathrm{fs}}^2 m_H/N_\phi,\quad N_\phi=13,\\
  m_p&\approx \alpha_{\mathrm{fs}}\, v/2\cdot(1+\kgeom^2/N_{12}),\\
  m_n&=m_p+2m_e\quad\text{(канал $\beta$-распада)}.
\end{align*}
\end{proposition}

Численное сравнение с данными PDG приведено в приложении; здесь зафиксированы формулы вывода.

\section{Четыре взаимодействия}

\begin{proposition}[Один закон --- четыре макроскопических канала]
Электромагнетизм: фазы на рёбрах, $\Phi_\square$, излучение с $n=0$.
Слабое: лево- и право-спиновые проекторы, переворот $SU(2)$; $\sin^2\theta_W=3/13$
из геометрии кластера $12+1$.
Сильное: составные узлы, $\alpha_s(v)=d/(N_{\mathrm{hier}}\pi)$.
Гравитация: деформация метрики на~$T$,
$G_{\mu\nu}=\mathcal{E}(\delta[\mathrm{strain}])=8\pi\ell_P^2 T_{\mu\nu}$.
\end{proposition}

Полный вывод гравитационных волн, полного спектра барионов и матрицы CKM
оставлен для дальнейшей работы и не ослабляет формулировок теорем выше.
